\section{迭代自我改进与反馈}

\subsection{Trace Feedback 与步级评分}
Chen 介绍 AlphaCode 的 trace feedback，将模型执行轨迹用自然语言 reify，形成 stepwise verifier。通过让模型自生成 line-by-line debug，它能识别哪些中间状态 promising，并用 text-level 评估影响后续 search。

\begin{figure}[H]
\centering
\includegraphics[width=0.92\textwidth]{frame-trace-feedback.jpg}
\caption{Trace feedback 示例画面，展示 line-by-line explanation 的可视化。\protect\footnotemark}
\label{fig:trace-feedback}
\end{figure}
\footnotetext{视频画面时间区间：01:11:20--01:12:15，讲解 trace feedback 的运作方式。}

\begin{importantbox}{Trace feedback 的高阶价值}
1. 把 step evaluation 转化为 natural language trace，使模型自己成为 verifier；2. AlphaCode 通过 trace 观察 context renovation，从而在 code generation 环节捕捉 subtle bugs；3. 更丰富的 trace 可直接作为 RL 奖励 signal，为 AlphaProof 与 AlphaGeometry 等系统提供 inference-time 改进路径。
\end{importantbox}

\subsection{自我修正提示与风险}
当没有 oracle verifier 时，自我修正往往把错误转换成另一种错误，Chen 以 01:13:23 处的负例提醒大家。设计 general-purpose feedback prompt 时，需突出 consistency、加入 trace context、并设定“不要翻动已确认部分”的 guard rails。

\begin{knowledgebox}{自我修正提示调参方向}
1. 强调 consistency：提示模型“列出两个可能的新步骤”，而非像“检查答案”这样模糊的引导；2. 注入 trace context：将 line-by-line explanation 再次作为 prompt，帮助模型用自身 reasoning 判断 correctness；3. 控制跳跃：增加“保留已确认推理”语句，防止模型因 feedback drift 把已验证的部分推翻。
\end{knowledgebox}

\begin{warningbox}{无 oracle 的 self-correction 坏处}
当只有 general feedback、没有 ground truth，模型会因为看不清哪条 response 可靠，反而保守地改写正确答案或放弃多步骤 reasoning，导致 accuracy 反而下降。
\end{warningbox}

\subsection{本章小结}
Trace feedback 与 self-correction 需要 reliable verifier 与 prompt guard rails，否则模型会在多个 round 后退步。设计 feedback 时必须兼顾 consistency 引导、trace context 和可解释性限制。

